本章では、量子回路シミュレータのパフォーマンスを測定し、NumPyの機能を最大限に活用して高速化する手法を学びます。プロファイリングによるボトルネックの特定、ペア単位のストライドアクセス、パウリゲートや2量子ビットゲートの構造を利用した特殊化、遅延規格化、そしてNumPy配列操作によるベクトル化までを段階的に実装します。最後に、回路実行時に最適な高速化パスを選択するゲートディスパッチの仕組みを構築します。

\begin{chaptergoal}
\begin{itemize}
\item プロファイリングで処理のボトルネックを特定し、ベンチマークで高速化の効果を測定できます。
\item 状態振幅へのストライドアクセスやゲートの構造を利用した特殊化、遅延規格化を実装できます。
\item NumPy配列操作によるベクトル化で、状態ベクトル更新を高速化できます。
\item 回路実行時に適切なゲート処理を選ぶディスパッチを構築できます。
\end{itemize}
\end{chaptergoal}

\section{パフォーマンス測定の基礎}

\subsection{なぜ高速化が必要か}

量子回路シミュレータの計算コストは、量子ビット数 $n$ に対して指数関数的に増大します。状態ベクトルは $2^n$ 個の複素数で表現されるため、量子ビット数が1増えるごとに必要なメモリ量は2倍になります。各ゲート適用には $O(2^n)$ の計算時間がかかるため、多数のゲートを含む回路では実行時間が膨大になります。

数学的には、1量子ビットゲートを全体の状態へ作用させる行列をクロネッカー積で書けます（\hyperref[sec:3-2-kron]{「クロネッカー積による全体の行列表現」}）。ただし、その全体行列は一般に $2^n\times2^n$ で、要素数は $4^n$ です。各量子ビットに対応する振幅ペアへ直接作用させれば、全体行列を保存せずに同じ計算を行えます。

量子ビット数の増加に伴って状態ベクトルのサイズと必要なメモリ量がどのように増大するかを表\ref{tab:memory_usage}に示します。

\begin{table}[htbp]
\centering
\begin{tabular}{crrl}
\hline
\textbf{量子ビット数} & \textbf{状態ベクトルサイズ} & \textbf{メモリ使用量} & \textbf{備考} \\ \hline
10 & 1,024 & 16 KB & 瞬時 \\
15 & 32,768 & 512 KB & 瞬時 \\
20 & 1,048,576 & 16 MB & 余裕 \\
25 & 33,554,432 & 512 MB & 現実的 \\
30 & 1,073,741,824 & 16 GB & ハイエンドPCで限界 \\
35 & 34,359,738,368 & 512 GB & スパコン級 \\ \hline
\end{tabular}
\caption{量子ビット数とメモリ使用量（complex128 型の場合）}
\label{tab:memory_usage}
\end{table}

この表から分かるように、30量子ビットを超えるシミュレーションには非常に大きなメモリが必要です。また、単にメモリに格納するだけでなく、すべての要素に対してゲート行列の演算を行う必要があるため、計算時間も同様に指数関数的に増加します。そこで、本書で実装してきた量子回路シミュレータを実用的な速度で動作させるために、さまざまな高速化手法を適用していきます。

高速化のアプローチは、大きく分けて以下の3つです。

\begin{itemize}
\item \textbf{アルゴリズム改善とNumPyベクトル化}: 状態ベクトル更新の計算量を削減します。ストライドアクセスによるペア単位の処理や、ゲート行列の構造を利用した特殊化、NumPyのスライスやブロードキャストによるベクトル化などが該当します（本章で解説）。
\item \textbf{JITコンパイル}: Numba を用いて Python のループを機械語にコンパイルし、実行速度を向上させます（第\ref{chap:11}章で解説）。
\item \textbf{GPUアクセラレーション}: CuPy を用いて GPU の多数のコアで並列計算を行います（第\ref{chap:11}章で解説）。
\end{itemize}

\section{プロファイリングによるボトルネック特定}

高速化の第一歩は、プログラムのどの部分が最も時間を消費しているかを特定することです。Python では \texttt{cProfile} モジュールを使ってプロファイリングを行えます。

\begin{lstlisting}[caption=cProfile を使ったプロファイリング（performance.py）, label=src:profiling]
import cProfile
import pstats
from quantum_sim import Circuit

def profile_circuit(num_qubits: int = 10, num_gates: int = 100):
    """回路実行のプロファイリングを行う。"""
    c = Circuit(num_qubits)
    for i in range(num_gates):
        target = i % num_qubits
        if i % 3 == 0:
            c.h(target)
        elif i % 3 == 1:
            c.x(target)
        else:
            c.cnot(target, (target + 1) % num_qubits)

    profiler = cProfile.Profile()
    profiler.enable()
    sv = c.statevector()
    profiler.disable()

    stats = pstats.Stats(profiler)
    stats.sort_stats("cumtime")
    stats.print_stats(15)  # 上位15件を表示

profile_circuit()
\end{lstlisting}

実行結果：

\begin{lstlisting}[caption=プロファイリング結果（抜粋）]
         1504 function calls in 0.018 seconds

   Ordered by: cumulative time

   ncalls  tottime  percall  cumtime  percall  filename:lineno
      100    0.010    0.000    0.015    0.000  qubit.py:76(apply_single)
       34    0.004    0.000    0.005    0.000  qubit.py:91(apply_two)
        1    0.001    0.001    0.018    0.018  circuit.py:165(statevector)
\end{lstlisting}

この結果から、\texttt{apply\_single} と \texttt{apply\_two} が全体の実行時間の大部分を占めていることが分かります。これらのメソッドが高速化の主要ターゲットです。

プロファイリング結果は \texttt{pstats} を使ってさまざまな観点で分析できます。\texttt{sort\_stats("cumtime")} は累積時間でソートしますが、\texttt{"time"} で内部実行時間のみ、\texttt{"ncalls"} で呼び出し回数でのソートも可能です。また、\texttt{print\_stats(15)} の数値を変えることで表示件数を調整できます。

\section{ベンチマーク測定ツール}

高速化の効果を定量的に評価するために、ウォームアップと複数回試行による平均実行時間を計測するベンチマーク用のクラス \texttt{PerformanceBenchmark} を実装します。\texttt{measure\_time} メソッドは回路を組み立てる関数を受け取り、組み立ては計測から外して \texttt{statevector()} の実行時間だけを測ります。

\begin{lstlisting}[caption=ベンチマーク計測クラス（performance.py、抜粋）, label=src:measure_time]
import time
from typing import Callable
import numpy as np
from quantum_sim import Circuit

class PerformanceBenchmark:
    """回路シミュレータのパフォーマンスベンチマークを行うクラス。"""

    def __init__(self, num_warmup: int = 5, num_runs: int = 20):
        self.num_warmup = num_warmup
        self.num_runs = num_runs

    def measure_time(
        self, circuit_fn: Callable[[], Circuit]
    ) -> BenchmarkResult:
        """回路の実行時間を計測する。"""
        # ウォームアップ
        for _ in range(self.num_warmup):
            c = circuit_fn()
            _ = c.statevector()

        # 本計測
        times = []
        for _ in range(self.num_runs):
            c = circuit_fn()
            t0 = time.perf_counter()
            _ = c.statevector()
            t1 = time.perf_counter()
            times.append((t1 - t0) * 1e6)  # マイクロ秒

        mean_t = float(np.mean(times))
        # 標準偏差・最小・最大・ゲート数も求め、計測結果をまとめた
        # データクラス BenchmarkResult にして返す（省略）
        ...
\end{lstlisting}

正確なベンチマークの測定では、以下の3点が重要です。

\begin{itemize}
\item \textbf{ウォームアップ}: 本測定の前に数回のウォームアップ実行を行い、JIT コンパイルやキャッシュの影響を安定化させます。
\item \textbf{複数回測定の平均}: システムの負荷変動や OS のスケジューリングによるノイズを抑えるため、複数回の平均値を採用します。
\item \texttt{\textbf{time.perf\_counter}}: マイクロ秒単位の微細な差を捉えるため、高分解能なタイマーを使用します。
\end{itemize}

\section{段階的な最適化の実装}

ここからは、量子回路シミュレータのコア部分である状態ベクトルへのゲート適用処理を、段階的に最適化していきます。本書のコードでは、最適化の各段階をメソッド名の接尾辞で区別しています。

\subsection{命名規則}

最適化の各段階を表すメソッド名の接尾辞と、その実装方針および対象読者の対応を表\ref{tab:naming_convention}にまとめます。

\begin{table}[htbp]
\centering
\begin{tabular}{cll}
\hline
\textbf{接尾辞} & \textbf{意味} & \textbf{対象読者} \\ \hline
\texttt{\_easy} & ビット演算を使わない教育的な実装 & 初心者 \\
\texttt{\_simple} & 最小限のビット演算による簡潔な実装 & 中級者 \\
\texttt{\_fast01} & 第一段階の高速化（ストライドアクセス in-place） & 上級者 \\
\texttt{\_fast02} & 第二段階の高速化（NumPy ベクトル化） & 上級者 \\
（なし） & \texttt{\_fast02} のエイリアス（最終最適化版） & 全員 \\ \hline
\end{tabular}
\caption{最適化段階の命名規則}
\label{tab:naming_convention}
\end{table}

この命名規則により、同じ機能を異なる手法で実装したメソッドを比較しながら学習できます。実際のコードでは、最終最適化版を接尾辞なしのメソッド名（\texttt{apply\_single} や \texttt{apply\_x} など）にエイリアスすることで、利用側は意識せずに最適化版を使えるようになっています。

\subsection{状態ベクトルの構造とゲート適用の考え方}

高速化を理解する上で重要なのは、$n$ 量子ビットの状態ベクトルにおいて、特定の量子ビットへのゲート適用がペア単位の操作になるという点です。

$n$ 量子ビットの状態ベクトル $\ket{\psi}$ は、$2^n$ 個のビット列基底 $\ket{i_{n-1} \dots i_1 i_0}$ の重ね合わせとして表されます。

\begin{equation}
\ket{\psi} = \sum_{i_0, i_1, \dots, i_{n-1} \in \{0,1\}}
\alpha_{i_{n-1} \dots i_1 i_0} \ket{i_{n-1} \dots i_1 i_0}
\end{equation}

ここで、標的となる量子ビット $k$ に1量子ビットゲート $G$ を適用する場合を考えます。量子ビット $k$ 以外のビットの値を固定すると、量子ビット $k$ が $0$ である成分と $1$ である成分のペアに対して、$G$ の $2 \times 2$ 行列が独立に作用します。

具体的には、量子ビット $k$ が $0$ である添字を $i$ とすると、量子ビット $k$ だけが $1$ である添字は $i + 2^k$ と表せます。この2つの振幅 $(\alpha_i, \alpha_{i + 2^k})^\top$ に対する $G$ の作用は次のようになります。

\begin{equation}
\begin{bmatrix}
\alpha'_i \\
\alpha'_{i + 2^k}
\end{bmatrix}
=
\begin{bmatrix}
G_{00} & G_{01} \\
G_{10} & G_{11}
\end{bmatrix}
\begin{bmatrix}
\alpha_i \\
\alpha_{i + 2^k}
\end{bmatrix}
\end{equation}

この操作は、状態ベクトルをストライド幅 $2^k$ で区切ったペア $(i, i + 2^k)$ ごとに $2 \times 2$ 行列を適用することに対応します。これがストライドアクセスの基本原理です。

状態ベクトルの成分がペアに分割され、各ペアに対して $2 \times 2$ ゲートが適用される対応関係を図\ref{fig:gate_pair_operation}に示します。

\begin{figure}[htbp]
\centering
\begin{tikzpicture}[
  font=\small,
  cell/.style={draw, minimum width=2.9cm, minimum height=0.52cm, align=center, inner sep=1pt},
  even/.style={fill=gray!12},
  odd/.style={fill=gray!28},
  gatebox/.style={draw, fill=white, rounded corners=2pt, minimum width=1.4cm, minimum height=0.88cm, align=center, font=\footnotesize},
  arrow/.style={-{Stealth[scale=0.9]}, thick}
]
  % 左側：ゲート適用前の状態ベクトル
  \node[font=\bfseries] at (1.45, 0.4) {状態ベクトル $\ket{\psi}$};
  \node[cell, even] (v0) at (1.45, -0.0) {$v[0]: \ket{00\mathbf{0}} \ (\alpha_{000})$};
  \node[cell, odd]  (v1) at (1.45, -0.55) {$v[1]: \ket{00\mathbf{1}} \ (\alpha_{001})$};
  \node[cell, even] (v2) at (1.45, -1.2) {$v[2]: \ket{01\mathbf{0}} \ (\alpha_{010})$};
  \node[cell, odd]  (v3) at (1.45, -1.75) {$v[3]: \ket{01\mathbf{1}} \ (\alpha_{011})$};
  \node[cell, even] (v4) at (1.45, -2.4) {$v[4]: \ket{10\mathbf{0}} \ (\alpha_{100})$};
  \node[cell, odd]  (v5) at (1.45, -2.95) {$v[5]: \ket{10\mathbf{1}} \ (\alpha_{101})$};
  \node[cell, even] (v6) at (1.45, -3.6) {$v[6]: \ket{11\mathbf{0}} \ (\alpha_{110})$};
  \node[cell, odd]  (v7) at (1.45, -4.15) {$v[7]: \ket{11\mathbf{1}} \ (\alpha_{111})$};

  % 中央：2x2 ゲート G
  \node[font=\footnotesize\bfseries, align=center, anchor=south] at (4.65, 0.28) {ペア単位の操作\\($2\times 2$ ゲート $G$)};
  \node[gatebox] (g0) at (4.65, -0.275) {$G = \begin{bmatrix} G_{00} & G_{01} \\ G_{10} & G_{11} \end{bmatrix}$};
  \node[gatebox] (g1) at (4.65, -1.475) {$G$};
  \node[gatebox] (g2) at (4.65, -2.675) {$G$};
  \node[gatebox] (g3) at (4.65, -3.875) {$G$};

  % 右側：ゲート適用後の状態ベクトル
  \node[font=\bfseries] at (7.85, 0.4) {適用後 $G\ket{\psi}$};
  \node[cell, even] (w0) at (7.85, -0.0) {$v'[0]: \ket{00\mathbf{0}} \ (\alpha'_{000})$};
  \node[cell, odd]  (w1) at (7.85, -0.55) {$v'[1]: \ket{00\mathbf{1}} \ (\alpha'_{001})$};
  \node[cell, even] (w2) at (7.85, -1.2) {$v'[2]: \ket{01\mathbf{0}} \ (\alpha'_{010})$};
  \node[cell, odd]  (w3) at (7.85, -1.75) {$v'[3]: \ket{01\mathbf{1}} \ (\alpha'_{011})$};
  \node[cell, even] (w4) at (7.85, -2.4) {$v'[4]: \ket{10\mathbf{0}} \ (\alpha'_{100})$};
  \node[cell, odd]  (w5) at (7.85, -2.95) {$v'[5]: \ket{10\mathbf{1}} \ (\alpha'_{101})$};
  \node[cell, even] (w6) at (7.85, -3.6) {$v'[6]: \ket{11\mathbf{0}} \ (\alpha'_{110})$};
  \node[cell, odd]  (w7) at (7.85, -4.15) {$v'[7]: \ket{11\mathbf{1}} \ (\alpha'_{111})$};

  % 接続線
  \draw[arrow] (v0.east) -- (g0.west |- v0);
  \draw[arrow] (v1.east) -- (g0.west |- v1);
  \draw[arrow] (g0.east |- w0) -- (w0.west);
  \draw[arrow] (g0.east |- w1) -- (w1.west);

  \draw[arrow] (v2.east) -- (g1.west |- v2);
  \draw[arrow] (v3.east) -- (g1.west |- v3);
  \draw[arrow] (g1.east |- w2) -- (w2.west);
  \draw[arrow] (g1.east |- w3) -- (w3.west);

  \draw[arrow] (v4.east) -- (g2.west |- v4);
  \draw[arrow] (v5.east) -- (g2.west |- v5);
  \draw[arrow] (g2.east |- w4) -- (w4.west);
  \draw[arrow] (g2.east |- w5) -- (w5.west);

  \draw[arrow] (v6.east) -- (g3.west |- v6);
  \draw[arrow] (v7.east) -- (g3.west |- v7);
  \draw[arrow] (g3.east |- w6) -- (w6.west);
  \draw[arrow] (g3.east |- w7) -- (w7.west);

  % 凡例
  \node[anchor=west, font=\footnotesize] at (0.2, -4.75) {※ $\text{target}=0$（最下位ビット）に着目};
  \node[draw, fill=gray!12, minimum width=0.4cm, minimum height=0.3cm, anchor=west] at (0.2, -5.25) {};
  \node[anchor=west, font=\footnotesize] at (0.75, -5.25) {ビット $=0$ (even)};
  \node[draw, fill=gray!28, minimum width=0.4cm, minimum height=0.3cm, anchor=west] at (3.4, -5.25) {};
  \node[anchor=west, font=\footnotesize] at (3.95, -5.25) {ビット $=1$ (odd)};
\end{tikzpicture}
\caption{状態ベクトルのペア分割とゲート適用の対応}
\label{fig:gate_pair_operation}
\end{figure}

\subsection{Step 1: ペア単位のストライドアクセス}

最初の最適化は、ゲート適用のループをペア単位のストライドアクセスに変更することです。

\texttt{apply\_single\_easy} は教育的な実装で、ビット演算を一切使わずにペアのインデックスを計算します。

\begin{lstlisting}[caption=apply\_single\_easy：教育版（ビット演算なし）（qubit.py）, label=src:apply_single_easy]
def apply_single_easy(self, gate: np.ndarray, target: int) -> None:
    """単一量子ビットゲート適用。教育版：ビット演算なし。"""
    stride = 2 ** target
    new_vec = np.zeros_like(self.vector)
    for block_start in range(0, self.size, stride * 2):
        for offset in range(stride):
            i = block_start + offset
            j = i + stride
            a = self.vector[i]
            b = self.vector[j]
            new_vec[i] = gate[0, 0] * a + gate[0, 1] * b
            new_vec[j] = gate[1, 0] * a + gate[1, 1] * b
    self.vector = new_vec
\end{lstlisting}

この実装では、\texttt{block\_start} がストライド幅の2倍ごとに進み、その内部で \texttt{offset} を使ってペアの前半分 $(i)$ と後半分 $(j = i + \text{stride})$ を直接計算します。ビット演算が一切登場しないため、ペア単位の操作が直感的に理解できます。

\texttt{apply\_single\_easy} からのさらなる改良として、ストライド幅の計算をビットシフト（\texttt{stride = 1 << target}）に変更した \texttt{apply\_single\_simple} や、\texttt{new\_vec} を介さず一時変数 \texttt{a, b} を用いて配列を直接 in-place に書き換える \texttt{apply\_single\_fast01} が挙げられます。
特に in-place 化により、ゲート適用ごとの \texttt{np.zeros\_like} によるメモリ確保が不要になり、キャッシュ局所性が向上します。これがストライドアクセスの基本形となります。

対象とする量子ビットの違いによってストライド幅と要素ペアの配置がどのように変化するかを図\ref{fig:stride_access}に示します。

\begin{figure}[htbp]
\centering
\begin{tikzpicture}[
  font=\small,
  cell/.style={draw, minimum width=1.15cm, minimum height=0.55cm, align=center, inner sep=1pt},
  even/.style={fill=gray!12},
  odd/.style={fill=gray!28},
  gate/.style={draw, fill=white, circle, inner sep=1.5pt, font=\scriptsize\bfseries},
  arcarrow/.style={thick, color=black}
]
  % ===== 上段: target = 0 (stride = 1) =====
  \node[anchor=west, font=\bfseries] at (0, 2.1) {$\text{target}=0$ （$\text{stride} = 2^0 = 1$、隣接ペア）};

  \foreach \i/\bin/\bg in {0/000/even, 1/001/odd, 2/010/even, 3/011/odd, 4/100/even, 5/101/odd, 6/110/even, 7/111/odd} {
    \node[cell, \bg] (c0_\i) at (1.25*\i + 0.65, 1.3) {$v[\i]$};
    \node[font=\tiny, text=gray!80!black] at (1.25*\i + 0.65, 1.75) {$\ket{\bin}$};
  }

  \foreach \i/\j in {0/1, 2/3, 4/5, 6/7} {
    \draw[arcarrow] (c0_\i.south) to[out=-60, in=-120] (c0_\j.south);
    \node[gate] at ($(c0_\i.south)!0.5!(c0_\j.south) + (0, -0.32)$) {$G$};
  }
  
  \draw[decorate, decoration={brace, amplitude=3pt, mirror}] 
    (c0_0.south west) -- (c0_0.south east) node[midway, below=3pt, font=\tiny] {$\text{stride}=1$};

  % ===== 下段: target = 1 (stride = 2) =====
  \node[anchor=west, font=\bfseries] at (0, -0.3) {$\text{target}=1$ （$\text{stride} = 2^1 = 2$、1つ飛ばしのペア）};

  \foreach \i/\bin/\bg in {0/000/even, 1/001/even, 2/010/odd, 3/011/odd, 4/100/even, 5/101/even, 6/110/odd, 7/111/odd} {
    \node[cell, \bg] (c1_\i) at (1.25*\i + 0.65, -1.1) {$v[\i]$};
    \node[font=\tiny, text=gray!80!black] at (1.25*\i + 0.65, -0.65) {$\ket{\bin}$};
  }

  \draw[arcarrow] (c1_0.south) to[out=-50, in=-130] (c1_2.south);
  \node[gate] at ($(c1_0.south)!0.5!(c1_2.south) + (0, -0.45)$) {$G$};

  \draw[arcarrow] (c1_1.south) to[out=-50, in=-130] (c1_3.south);
  \node[gate] at ($(c1_1.south)!0.5!(c1_3.south) + (0, -0.45)$) {$G$};

  \draw[arcarrow] (c1_4.south) to[out=-50, in=-130] (c1_6.south);
  \node[gate] at ($(c1_4.south)!0.5!(c1_6.south) + (0, -0.45)$) {$G$};

  \draw[arcarrow] (c1_5.south) to[out=-50, in=-130] (c1_7.south);
  \node[gate] at ($(c1_5.south)!0.5!(c1_7.south) + (0, -0.45)$) {$G$};

  \draw[decorate, decoration={brace, amplitude=4pt, mirror}] 
    (c1_0.south west) -- (c1_1.south east) node[midway, below=3pt, font=\tiny] {$\text{stride}=2$};

  % 凡例
  \node[draw, fill=gray!12, minimum width=0.35cm, minimum height=0.25cm, anchor=west] at (2.5, -2.3) {};
  \node[anchor=west, font=\footnotesize] at (2.9, -2.3) {targetビット $=0$};
  \node[draw, fill=gray!28, minimum width=0.35cm, minimum height=0.25cm, anchor=west] at (5.8, -2.3) {};
  \node[anchor=west, font=\footnotesize] at (6.2, -2.3) {targetビット $=1$};
\end{tikzpicture}
\caption{ストライドアクセスの概念図}
\label{fig:stride_access}
\end{figure}

\subsection{Step 2: パウリゲートの特殊化}

パウリゲート $X, Y, Z$ は非常に単純な行列構造を持つため、一般の2x2行列積として計算するのは非効率です。各ゲートの構造を利用した専用メソッドを実装することで、複素乗算を削減できます。

各パウリゲートの行列表現と、それらに対応する簡略化されたペア操作の内容を表\ref{tab:pauli_specialization}にまとめます。

\begin{table}[htbp]
\centering
\begin{tabular}{cll}
\hline
\textbf{ゲート} & \textbf{行列} & \textbf{ペア操作} \\ \hline
$X$ & $\begin{bmatrix}0&1\\1&0\end{bmatrix}$ & \texttt{a, b = b, a}（値の入れ替えのみ） \\
$Y$ & $\begin{bmatrix}0&-\ii\\ \ii&0\end{bmatrix}$ & \texttt{a, b = $-\ii b$, $\ii a$}（乗算2回のみ） \\
$Z$ & $\begin{bmatrix}1&0\\0&-1\end{bmatrix}$ & \texttt{b *= -1}（符号反転のみ） \\
$H$ & $\frac{1}{\sqrt{2}}\begin{bmatrix}1&1\\1&-1\end{bmatrix}$ & 一般2x2行列適用（Step 1の方式） \\ \hline
\end{tabular}
\caption{パウリゲートの特殊化}
\label{tab:pauli_specialization}
\end{table}

\textbf{X ゲート}は、ペア $(a, b)$ を $(b, a)$ に入れ替えるだけです。複素乗算は一切不要で、値のスワップだけで済みます。

\begin{lstlisting}[caption=apply\_x：パウリXゲート特殊化（qubit.py）, label=src:apply_x]
def apply_x_fast01(self, target: int) -> None:
    """パウリ X ゲート特殊化。
    X = [[0,1],[1,0]] → ペア (a,b) を (b,a) に入れ替え。"""
    stride = 1 << target
    v = self.vector
    for i in range(0, self.size, stride * 2):
        for j in range(i, i + stride):
            v[j], v[j + stride] = v[j + stride], v[j]
\end{lstlisting}

\textbf{Z ゲート}は、target ビットが1の要素（ペアの後半）に $-1$ を乗算するだけです。

\begin{lstlisting}[caption=apply\_z：パウリZゲート特殊化（qubit.py）, label=src:apply_z]
def apply_z_fast01(self, target: int) -> None:
    """パウリ Z ゲート特殊化。
    Z = [[1,0],[0,-1]] → target ビットが1の要素に -1 を乗算。"""
    stride = 1 << target
    v = self.vector
    for i in range(0, self.size, stride * 2):
        for j in range(i + stride, i + stride * 2):
            v[j] *= -1
\end{lstlisting}

$Y$ ゲート（乗算2回）や $H$ ゲートも表\ref{tab:pauli_specialization}に示したとおり同様のペア操作として効率化できます。このように一般の行列積をゲート固有の軽量な操作に置き換えることで、不要な複素乗算を徹底的に排除できます。

\subsection{Step 3: CNOT/CZ/SWAP の特殊化}

2量子ビットゲートも、パウリゲートと同様に構造を利用した特殊化が可能です。これらのゲートは4x4行列として実装できますが、その構造を利用すれば要素のスワップや符号反転だけで実装できます。

\textbf{CNOT ゲート}は、制御ビットが1のとき標的ビットを反転（X適用）します。4要素ブロックのうち制御ビット=1に対応する2要素をスワップします。

\begin{lstlisting}[caption=apply\_cnot：CNOTゲート特殊化（qubit.py）, label=src:apply_cnot]
def apply_cnot_fast01(self, control: int, target: int) -> None:
    """CNOT ゲート特殊化。
    制御ビット=1 のとき標的ビットを反転（X適用）。
    4要素ブロックのうち制御=1 に対応する2要素をスワップ。"""
    s1, s2 = 1 << control, 1 << target
    step = 2 * max(s1, s2)
    sub_step = 2 * min(s1, s2)
    sub_size = min(s1, s2)
    v = self.vector
    for outer in range(0, self.size, step):
        for mid in range(outer, outer + max(s1, s2), sub_step):
            for off in range(mid, mid + sub_size):
                i10 = off + s1        # control=1, target=0
                i11 = off + s1 + s2   # control=1, target=1
                v[i10], v[i11] = v[i11], v[i10]
\end{lstlisting}

\textbf{CZ ゲート}や \textbf{SWAP ゲート}も、最内層ループの操作が異なるだけです。CZ では両ビットが1の要素のみ符号反転（\texttt{v[off + s1 + s2] *= -1}）し、SWAP では $|01\rangle$ と $|10\rangle$ に対応する要素（\texttt{i10 = off + s1} と \texttt{i01 = off + s2}）を入れ替えます。いずれも一般の 4x4 行列積を回避し、最小限のメモリ操作で完了します。

\subsection{Step 4: 遅延規格化}

ユニタリゲートはノルムを保存するため、ゲート適用のたびに \texttt{\_normalize()} を呼ぶ必要はありません。\texttt{StateVector} クラスでは、\texttt{\_needs\_normalize} フラグを使って規格化を遅延実行しています。

\begin{lstlisting}[caption=遅延規格化の実装（qubit.py）, label=src:delayed_normalize]
class StateVector:
    def __init__(self, num_qubits: int = 1) -> None:
        # ...（省略）...
        self._needs_normalize: bool = False

    def _normalize(self) -> None:
        """状態ベクトルを規格化する。"""
        norm = np.linalg.norm(self.vector)
        if norm > 1e-15:
            self.vector /= norm
        self._needs_normalize = False

    def _ensure_normalized(self) -> None:
        """必要に応じて規格化する（遅延規格化）。"""
        if self._needs_normalize:
            self._normalize()

    def probabilities(self) -> np.ndarray:
        self._ensure_normalized()
        return np.abs(self.vector) ** 2

    def measure(self, qubit: int, rng=None) -> int:
        self._ensure_normalized()
        # ...（測定処理）...
\end{lstlisting}

ユニタリゲートの適用後はノルムが保存されるため、\texttt{\_needs\_normalize} は \texttt{False} のままです。一方、測定による射影など非ユニタリ操作の後ではフラグが \texttt{True} に設定され、次回の \texttt{probabilities()} や \texttt{measure()} 呼び出し時に自動的に規格化が実行されます。この仕組みにより、不要なノルム計算を省くことができます。

\subsubsection{遅延規格化の効果}

10量子ビット・100ゲートの回路で測定したところ、各ゲート後に逐次規格化した場合（799.5μs）に対し、測定直前まで遅延させた場合は 496.7μs と約 \textbf{1.6倍（38\%短縮）} の高速化が確認できました。
\texttt{\_normalize} 1回あたり約 2.8μs、\texttt{apply\_h\_fast02} 1回あたり約 4.8μs を要するため、ゲート適用のたびに規格化するとゲート本体の半分以上の時間を規格化に浪費することになります。遅延規格化はこの無駄を省き、特に高速化されたゲート実行において安定した高速化をもたらします。

\subsection{Step 5: NumPy ベクトル化（単一量子ビットゲート）}

Python の for ループは、C言語などと比較すると実行速度が遅いという欠点があります。NumPy のベクトル化演算を活用すれば、Python ループを排除し、NumPy の内部 C 実装に処理を委譲できます。

\texttt{apply\_single\_fast02} は、reshape + スライスでペアを分割し、NumPy のブロードキャストで一括演算を行います。

\begin{lstlisting}[caption=apply\_single\_fast02：NumPy ベクトル化（qubit.py）, label=src:apply_single_fast02]
def apply_single_fast02(self, gate: np.ndarray, target: int) -> None:
    """単一量子ビットゲート適用。高速化②：numpy ベクトル化。
    reshape + スライスでペアを分割し、numpy のブロードキャストで一括演算。"""
    stride = 1 << target
    v = self.vector
    # ペアに分割: even = targetビットが0, odd = targetビットが1
    even = v.reshape(-1, 2 * stride)[:, :stride].ravel()
    odd = v.reshape(-1, 2 * stride)[:, stride:2 * stride].ravel()
    # 行列演算をベクトル化
    new_even = gate[0, 0] * even + gate[0, 1] * odd
    new_odd = gate[1, 0] * even + gate[1, 1] * odd
    # 書き戻し
    v.reshape(-1, 2 * stride)[:, :stride] = new_even.reshape(-1, stride)
    v.reshape(-1, 2 * stride)[:, stride:2 * stride] = new_odd.reshape(-1, stride)
\end{lstlisting}

このコードの動作を理解するために、具体例で考えてみましょう。target=0（最下位ビット）の場合、stride=1 です。状態ベクトルを \texttt{reshape(-1, 2)} で2列の行列に変形すると、0列目が偶数インデックス（targetビット=0）、1列目が奇数インデックス（targetビット=1）に対応します。これを行列 $G$ で変換した後、元の形状に戻します。

target=1 の場合は stride=2 となり、\texttt{reshape(-1, 4)} で4列の行列に変形します。このうち最初の2列が targetビット=0、残りの2列が targetビット=1 に対応します。このように、\texttt{reshape} とスライスを組み合わせることで、任意の target に対するペア分割をベクトル化して表現できます。

この reshape とスライスによる状態ベクトルの変形とペア分割、および一括演算の手順を図\ref{fig:numpy_vectorization}に示します。

\begin{figure}[htbp]
\centering
\begin{tikzpicture}[
  font=\small,
  cell/.style={draw, minimum width=1.15cm, minimum height=0.55cm, align=center, inner sep=1pt},
  even/.style={fill=gray!12},
  odd/.style={fill=gray!28},
  arrow/.style={-{Stealth[scale=0.9]}, thick},
  codebox/.style={draw, dashed, fill=gray!5, rounded corners=2pt, font=\footnotesize\ttfamily, align=left, inner sep=4pt}
]
  % ===== 1段目: 1次元配列 v =====
  \node[anchor=west, font=\bfseries] at (0, 4.3) {① 1次元状態ベクトル \texttt{v}（shape: \texttt{(8,)}、$\text{target}=1, \text{stride}=2$ の例）};
  
  \foreach \i/\bg in {0/even, 1/even, 2/odd, 3/odd, 4/even, 5/even, 6/odd, 7/odd} {
    \node[cell, \bg] (v\i) at (1.25*\i + 0.65, 3.6) {$v[\i]$};
  }

  % 矢印 1 -> 2
  \draw[arrow] (5.0, 3.15) -- (5.0, 2.35) 
    node[midway, right=4pt, font=\footnotesize\ttfamily] {v.reshape(-1, 2 * stride) \rmfamily （ビュー変換）};

  % ===== 2段目: 2次元配列 =====
  \node[anchor=west, font=\bfseries] at (0, 2.65) {② 2次元配列に変形};
  \node[anchor=west, font=\footnotesize] at (7.15, 1.05) {shape: \texttt{(2, 4)}};

  \node[cell, even] (m00) at (2.6, 1.35) {$v[0]$};
  \node[cell, even] (m01) at (3.85, 1.35) {$v[1]$};
  \node[cell, odd]  (m02) at (5.1, 1.35) {$v[2]$};
  \node[cell, odd]  (m03) at (6.35, 1.35) {$v[3]$};

  \node[cell, even] (m10) at (2.6, 0.75) {$v[4]$};
  \node[cell, even] (m11) at (3.85, 0.75) {$v[5]$};
  \node[cell, odd]  (m12) at (5.1, 0.75) {$v[6]$};
  \node[cell, odd]  (m13) at (6.35, 0.75) {$v[7]$};

  % スライスのブレース
  % 見出しと重ならないよう、ブレースは配列の右側に縦向きで置く
  \draw[decorate, decoration={brace, amplitude=4pt}]
    (2.0, 1.72) -- (4.45, 1.72)
    node[midway, above=3pt, font=\scriptsize\ttfamily] {[:, :stride]};
  \draw[decorate, decoration={brace, amplitude=4pt}]
    (4.5, 1.72) -- (6.95, 1.72)
    node[midway, above=3pt, font=\scriptsize\ttfamily] {[:, stride:2*stride]};

  % 矢印 2 -> 3
  \draw[arrow] (3.2, 0.35) -- (2.6, -0.35);
  \draw[arrow] (5.7, 0.35) -- (6.3, -0.35);

  % ===== 3段目: スライス抽出と一括ベクトル演算 =====
  \node[anchor=west, font=\bfseries] at (0, -0.55) {③ 前半・後半スライスの切り出しと一括行列演算};

  % even 配列
  \node[anchor=west, font=\footnotesize\ttfamily] at (0.3, -0.95) {even \rmfamily （\ttfamily [:, :2]\rmfamily ）};
  \node[cell, even] at (0.9, -1.45) {$v[0]$};
  \node[cell, even] at (2.1, -1.45) {$v[1]$};
  \node[cell, even] at (3.3, -1.45) {$v[4]$};
  \node[cell, even] at (4.5, -1.45) {$v[5]$};

  % odd 配列
  \node[anchor=west, font=\footnotesize\ttfamily] at (5.5, -0.95) {odd \rmfamily （\ttfamily [:, 2:4]\rmfamily ）};
  \node[cell, odd] at (6.1, -1.45) {$v[2]$};
  \node[cell, odd] at (7.3, -1.45) {$v[3]$};
  \node[cell, odd] at (8.5, -1.45) {$v[6]$};
  \node[cell, odd] at (9.7, -1.45) {$v[7]$};

  % 下部コード注記
  \node[codebox, anchor=north, minimum width=10.5cm] at (5.0, -1.95) {
    \textbf{\rmfamily ベクトル化演算（NumPy ブロードキャスト）:}\\
    \texttt{new\_even = gate[0,0] * even + gate[0,1] * odd}\\
    \texttt{new\_odd  = gate[1,0] * even + gate[1,1] * odd}\\
    \texttt{v.reshape(-1, 4)[:, :2] = new\_even.reshape(-1, 2)  \# 書き戻し}\\
    \texttt{v.reshape(-1, 4)[:, 2:4] = new\_odd.reshape(-1, 2)}
  };
\end{tikzpicture}
\caption{NumPy ベクトル化の模式図（reshape + スライス）}
\label{fig:numpy_vectorization}
\end{figure}

パウリゲートの特殊化にも、同様のベクトル化版が用意されています。X ゲートでは、スライスで抽出した \texttt{even} と \texttt{odd} を入れ替えて書き戻します。

\begin{lstlisting}[caption=パウリXゲート特殊化のベクトル化版（qubit.py）, label=src:pauli_vectorized]
def apply_x_fast02(self, target: int) -> None:
    """パウリ X ゲート特殊化。numpy ベクトル化。"""
    stride = 1 << target
    v = self.vector
    even = v.reshape(-1, 2 * stride)[:, :stride].copy()
    odd = v.reshape(-1, 2 * stride)[:, stride:2 * stride].copy()
    v.reshape(-1, 2 * stride)[:, :stride] = odd.reshape(-1, stride)
    v.reshape(-1, 2 * stride)[:, stride:2 * stride] = even.reshape(-1, stride)
\end{lstlisting}

なお Z ゲートの場合は、\texttt{odd = v.reshape(-1, 2 * stride)[:, stride:2 * stride]} に対して \texttt{odd *= -1} とするだけで直接符号反転でき、配列コピーすら不要です。

\subsection{Step 6: 2量子ビットゲートのベクトル化}

2量子ビットゲートも同様にベクトル化できます。\texttt{apply\_two\_fast02} は、4要素ブロックのインデックス配列を生成し、一括で行列演算を行います。

\begin{lstlisting}[caption=\_two\_block\_indices：4要素ブロックのインデックス生成（qubit.py）, label=src:two_block_indices]
@staticmethod
def _two_block_indices(target0: int, target1: int, size: int):
    """2量子ビットゲートの4要素ブロックインデックスを生成する。
    4要素ブロック (i00, i10, i01, i11) を正しく一度ずつ巡回する。"""
    s1, s2 = 1 << target0, 1 << target1
    step = 2 * max(s1, s2)
    sub_step = 2 * min(s1, s2)
    sub_size = min(s1, s2)
    n_outer = size // step
    n_inner = max(s1, s2) // sub_step

    base = np.arange(n_outer)[:, np.newaxis, np.newaxis] * step
    mid = np.arange(n_inner)[np.newaxis, :, np.newaxis] * sub_step
    off = np.arange(sub_size)[np.newaxis, np.newaxis, :]

    k = (base + mid + off).ravel()
    return k, k + s1, k + s2, k + s1 + s2
\end{lstlisting}

この \texttt{\_two\_block\_indices} メソッドは、NumPy のブロードキャストを活用して4要素ブロック $(i_{00}, i_{10}, i_{01}, i_{11})$ のインデックス配列を生成します。このインデックス配列を使うことで、Python の三重ループを排除できます。

\begin{lstlisting}[caption=apply\_two\_fast02：2量子ビットゲートのベクトル化（qubit.py）, label=src:apply_two_fast02]
def apply_two_fast02(self, gate: np.ndarray, target0: int, target1: int) -> None:
    """2量子ビットゲート適用。高速化②：numpy ベクトル化。"""
    i00, i10, i01, i11 = self._two_block_indices(target0, target1, self.size)
    v = self.vector
    a00 = v[i00]
    a10 = v[i10]
    a01 = v[i01]
    a11 = v[i11]
    v[i00] = gate[0, 0] * a00 + gate[0, 1] * a10 + gate[0, 2] * a01 + gate[0, 3] * a11
    v[i10] = gate[1, 0] * a00 + gate[1, 1] * a10 + gate[1, 2] * a01 + gate[1, 3] * a11
    v[i01] = gate[2, 0] * a00 + gate[2, 1] * a10 + gate[2, 2] * a01 + gate[2, 3] * a11
    v[i11] = gate[3, 0] * a00 + gate[3, 1] * a10 + gate[3, 2] * a01 + gate[3, 3] * a11
\end{lstlisting}

2量子ビットゲートの特殊化（CNOT, CZ, SWAP）にも、同様にベクトル化版を定義できます。代表例として CNOT では、\texttt{\_two\_block\_indices} で得たインデックスを用いて $i_{10}$ と $i_{11}$ の要素を一括スワップします。

\begin{lstlisting}[caption=CNOT 特殊化のベクトル化版（qubit.py）, label=src:twoq_vectorized]
def apply_cnot_fast02(self, control: int, target: int) -> None:
    """CNOT ゲート特殊化。numpy ベクトル化。"""
    _, i10, _, i11 = self._two_block_indices(control, target, self.size)
    v = self.vector
    tmp = v[i10].copy()
    v[i10] = v[i11]
    v[i11] = tmp
\end{lstlisting}

CZ ゲートの場合は $i_{11}$ の要素を一括で符号反転（\texttt{self.vector[i11] *= -1}）し、SWAP ゲートの場合は $i_{10}$ と $i_{01}$ を一括スワップするだけで実装できます。

\subsection{Step 7: ゲート融合（回路レベル）}

ゲート融合は、第\ref{chap:9}章で実装した \texttt{fuse\_adjacent\_gates} と連携する最適化です。隣接する同一量子ビットへのゲートを事前に行列積で融合することで、ゲート適用回数を削減します。

例えば、$H(q_0) \to X(q_0) \to Z(q_0)$ という3つの連続するゲートは、事前に $Z \cdot X \cdot H$ の行列積を計算しておき、1回のゲート適用で済ませることができます。

\begin{equation}
Z \cdot X \cdot H = \begin{bmatrix}1&0\\0&-1\end{bmatrix}
\begin{bmatrix}0&1\\1&0\end{bmatrix}
\frac{1}{\sqrt{2}}\begin{bmatrix}1&1\\1&-1\end{bmatrix}
= \frac{1}{\sqrt{2}}\begin{bmatrix}1&-1\\-1&-1\end{bmatrix}
\end{equation}

この最適化は回路の深さ（ゲート数）が大きいほど効果的です。第\ref{chap:9}章の \texttt{CircuitOptimizer} でレベル3以上を設定すると、ゲート融合が有効になります。

本章で解説してきた状態ベクトル更新の高速化手法と独立した各種最適化手法の全体像を図\ref{fig:optimization_map}に示します。

\begin{figure}[htbp]
\centering
\begin{tikzpicture}[
  font=\small,
  node distance=0.6cm and 0.4cm,
  stepbox/.style={draw, rounded corners=3pt, fill=white, line width=0.8pt, align=center, inner sep=4pt},
  mainstep/.style={stepbox, fill=gray!8},
  indepbox/.style={stepbox, fill=gray!3, dashed},
  arrow/.style={-{Stealth[scale=0.9]}, thick},
  labeltext/.style={font=\scriptsize, text=gray!80!black}
]
  % ===== メイン系統: 状態ベクトル更新の最適化 =====
  \node[mainstep, minimum width=7.2cm] (step1) at (0, 4.0) {
    \textbf{Step 1: ペア単位のストライドアクセス}\\
    \footnotesize 状態ベクトルを $2^{\text{target}}$ 刻みで巡回・in-place 更新
  };

  \node[mainstep, below left=0.75cm and -2.4cm of step1, minimum width=3.4cm] (step2) {
    \textbf{Step 2: パウリゲートの特殊化}\\
    \footnotesize $X$: 入替, $Z$: 符号反転\\
    乗算を削減・排除
  };

  \node[mainstep, below right=0.75cm and -2.4cm of step1, minimum width=3.4cm] (step3) {
    \textbf{Step 3: CNOT/CZ/SWAP の特殊化}\\
    \footnotesize 4要素ブロック直操作\\
    一般4x4行列積を回避
  };

  \node[mainstep, below=0.75cm of step2, minimum width=3.4cm] (step5) {
    \textbf{Step 5: NumPy ベクトル化}\\
    \footnotesize （単一量子ビット）\\
    reshape + スライス一括演算
  };

  \node[mainstep, below=0.75cm of step3, minimum width=3.4cm] (step6) {
    \textbf{Step 6: 2量子ビットゲートの}\\
    \textbf{ベクトル化}\\
    \footnotesize ブロックインデックス一括演算
  };

  % 矢印
  \draw[arrow] (step1.south -| step2.north) -- (step2.north) 
    node[midway, left, labeltext] {演算削減};
  \draw[arrow] (step1.south -| step3.north) -- (step3.north) 
    node[midway, right, labeltext] {演算削減};

  \draw[arrow] (step2.south) -- (step5.north) 
    node[midway, left, labeltext] {ループ排除};
  \draw[arrow] (step3.south) -- (step6.north) 
    node[midway, right, labeltext] {ループ排除};

  % ===== 独立した最適化 =====
  \node[indepbox, minimum width=4.3cm] (step4) at (-2.6, -2.75) {
    \textbf{Step 4: 遅延規格化}\\
    \footnotesize ユニタリ後のノルム計算省略\\
    測定・確率取得時のみ実行
  };

  \node[indepbox, minimum width=4.3cm] (step7) at (2.6, -2.75) {
    \textbf{Step 7: ゲート融合（回路レベル）}\\
    \footnotesize 連続ゲートを行列積で事前合成\\
    ゲート適用回数を削減
  };

  % グループ枠
  \node[draw=gray!60, dotted, rounded corners=4pt, fit=(step4) (step7), inner sep=6pt, label={[font=\footnotesize\bfseries, text=gray!80!black]below:独立した最適化（メイン系統と直交して効く）}] (indepgroup) {};
  \node[draw=gray!60, dotted, rounded corners=4pt, fit=(step1) (step2) (step3) (step5) (step6), inner sep=6pt, label={[font=\footnotesize\bfseries, text=gray!80!black]above:状態ベクトル更新の高速化}] (maingroup) {};
\end{tikzpicture}
\caption{最適化の全体マップ}
\label{fig:optimization_map}
\end{figure}

\section{Circuit のゲートディスパッチ}

本章で実装した各種の特殊化メソッドは、第\ref{chap:6}章のリスト\ref{src:circuit_class}で示した \texttt{Circuit} クラスの \texttt{\_apply\_gate} メソッドにおいてゲート名に応じて呼び分けられます。X, Y, Z, H, CNOT, CZ, SWAP などの既知のゲートは複素乗算を省いた高速な専用パスへ割り振られ、回転ゲートや一般のユニタリゲートのみが汎用行列演算パスを通ることで、回路全体の実行速度が大幅に向上します。

\section{この章で学んだこと}

\begin{itemize}
\item 量子回路シミュレータの計算コストは量子ビット数に対して指数関数的に増大するため、ボトルネック箇所のプロファイリングと段階的な最適化が不可欠です。
\item \textbf{ストライドアクセス}により、状態ベクトルの対象ビットが0と1のペアを直接インデックス計算して in-place 更新できます。
\item パウリゲート（X, Y, Z）や基本ゲート（CNOT, CZ, SWAP）の代数的性質を利用し、複素行列積を要素の入れ替えや符号反転に置き換える\textbf{ゲート特殊化}が極めて有効です。
\item \textbf{遅延規格化}により、毎ステップのノルム計算を省き測定直前のみに限定することで、不要な割算コストを削減できます。
\item \textbf{NumPyベクトル化}により、Pythonの \texttt{for} ループを排除してC言語レベルの一括スライス操作に移行することで、大幅な高速化を達成できます。
\item \texttt{Circuit} に\textbf{ゲートディスパッチ}を組み込むことで、既知のゲートには特殊化・ベクトル化メソッドを、一般ゲートには汎用行列積を自動的に割り振ることができます。
\end{itemize}

NumPyによるアルゴリズム改善とベクトル化によってCPU上の実行速度は大きく向上しましたが、次章（第\ref{chap:11}章）では、さらなる極限の性能を求めて、NumbaによるJITコンパイルやCuPyによるGPUアクセラレーション、マルチスレッド並列サンプリングを導入します。
